\section{财产犯罪 - 罪数联系}
\begin{enumerate}
  \item 大总结，综合运用。\bs{核心。先找被害人}。
  \item 销赃行为。
        \begin{enumerate}
          \item 盗窃 + 销赃。\al{甲对乙指着丙的树说这是我的，乙交其 3 万元，乙便伐倒运走。盗窃、诈骗择一重}。
          \item 诈骗 + 销赃。两头骗。\al{甲骗到乙的车，然后向丙谎称是自己的，卖给了丙。诈骗 + 诈骗并罚}。
          \item 侵占 + 销赃。\al{甲租用乙汽车，到期拒不归还，又谎称是自己所有卖给丙。侵占、诈骗择一重}。
        \end{enumerate}
  \item 免除返还义务。
        \begin{enumerate}
          \item 先侵占后诈骗。\al{甲代为保管乙财物，自行变卖，对乙谎称被偷。侵占、诈骗}。争议。是否并罚。
                \begin{enumerate}
                  \item \zd{多数说}。重罪吸收轻罪，定诈骗。
                  \item 少数说。只定侵占，诈骗行为不可罚的事后行为。
                \end{enumerate}
          \item \bs{先侵占后杀害}。\al{金池长老借袈裟 \yjt 故意杀人 + 抢劫致人死亡（故意）\yjt 想象竞合择一重}。
                \begin{enumerate}
                  \item 杀人重 \yjt 杀人 + 侵占数罪并罚。
                  \item 抢劫致人死亡（重）\yjt 抢劫致人死亡 + 侵占 \yjt \zd{重罪吸收轻罪} \yjt 抢劫致人死亡。
                \end{enumerate}
          \item 先诈骗后抢劫。\al{甲骗到乙的画，乙发现后要求返还，甲暴力威胁。诈骗、抢劫}。是否并罚同上。
        \end{enumerate}
  \item 先偷后骗。\al{甲在超市偷了一部手机，第二天去退货退款。盗窃、诈骗}。
  \item 后续侵占。
        \begin{enumerate}
          \item 侵占后侵占（连续侵占）。\al{甲保管乙电脑，乙 3 日后请求返还，甲拒不返还并出卖给丙}。出卖行为无权处分，依然是侵占。连续犯。
          \item 取得后侵占。
        \end{enumerate}
  \item \bs{冒名实现债权}。\yi \bs{多数说。构成三角诈骗}。\er 实务界。定盗窃。
        \begin{enumerate}
          \item \al{偷换二维码}。
          \item \al{甲借用乙的账号网购，甲瞒着乙联系商家谎称收件地址错误，改为自己的地址。商家将货物寄给了甲}。
        \end{enumerate}
\end{enumerate}
